\documentclass[../../../include/open-logic-section]{subfiles}

\begin{document}

\olfileid[hi]{sth}{z}{infinity-again}	
\olsection{अनंतता}

हमारे पास यह सुनिश्चित करने के लिए पर्याप्त अभिगृहीत पहले से हैं कि अनंततः
अनेक समुच्चय हैं (यदि कोई समुच्चय है)। मान लें कि किसी समुच्चय का अस्तित्व
है, अतः $\emptyset$ का अस्तित्व है
(\olref[sth][z][sep]{prop:emptyexists} से)। अब किसी भी समुच्चय~$x$ के लिए
\olref[sth][z][pairs]{prop:pairsconsequences} से समुच्चय $x \cup \{x\}$
का अस्तित्व है। इसे कुछ बार लागू करने पर हमें निम्न समुच्चय मिलेंगे:
\begin{enumerate}
	\item[0.] $\emptyset$ %& $0$ & stage $0$\\
	\item[1.] $ \{\emptyset\}$ %& $1$ & stage $1$\\
	\item[2.] $\{\emptyset, \{\emptyset\}\}$ %& $2$ & stage $2$\\
	\item[3.] $\{\emptyset, \{\emptyset\}, \{\emptyset, \{\emptyset\}\}\}$ %&$3$ & stage $3$\\
	\item[4.] $\{\emptyset, \{\emptyset\}, \{\emptyset, \{\emptyset\}\}, \{\emptyset, \{\emptyset\}, \{\emptyset, \{\emptyset\}\}\}\}$% & $4$ & stage $4$ 
\end{enumerate}
और हम जाँच सकते हैं कि इनमें से प्रत्येक समुच्चय भिन्न है।

हमने कुछ कारणों से क्रमांकन $0$ से आरंभ किया है। उनमें से एक यह है। यह
जाँचना बहुत कठिन नहीं कि ``$n$'' अंकित समुच्चय के ठीक $n$ सदस्य हैं और
(सहजतः) वह $n$वें चरण पर बनता है।

किंतु इससे हमें \emph{अनंततः अनेक} समुच्चय तो मिलते हैं, यह गारंटी नहीं
मिलती कि कोई \emph{अनंत समुच्चय}, अर्थात अनंततः अनेक सदस्यों वाला समुच्चय,
है। और यह सचमुच महत्त्वपूर्ण है: जब तक कोई (डेडेकिंड-)अनंत समुच्चय न मिले,
हम डेडेकिंड बीजगणित नहीं बना सकते। पर हमें डेडेकिंड बीजगणित चाहिए ताकि उसे
प्राकृतिक संख्याओं का समुच्चय मान सकें।
(\olref[sfr][infinite][dedekindsproof]{sec} से तुलना करें।)

महत्त्वपूर्ण बात यह है कि अब तक निर्धारित अभिगृहीत किसी अनंत समुच्चय के
अस्तित्व की गारंटी \emph{नहीं} देते। अतः हमें नया अभिगृहीत निर्धारित करना
होगा:

\begin{axiom}[अनंतता]
ऐसा समुच्चय $I$ है कि $\emptyset \in I$, और जब भी $x \in I$ हो तब
$x \cup \{x\} \in I$।
\begin{align*}
	\exists I( & (\exists o \in I)\forall x\ x \notin o \land {}\\
	& (\forall x \in I)(\exists s \in I)\forall z(z \in s \liff (z \in x \lor z = x)))
\end{align*}
\end{axiom}

यह देखना सरल है कि अनंतता के अभिगृहीत से मिला समुच्चय $I$ डेडेकिंड-अनंत
है। उसका विशिष्ट अवयव $\emptyset$ है और $I$ पर अंतःक्षेप
$s(x) = x\cup \{x\}$ से मिलता है। अब
\olref[sfr][infinite][dedekind]{thm:DedekindInfiniteAlgebra} ने दिखाया कि
किसी डेडेकिंड-अनंत समुच्चय से डेडेकिंड बीजगणित कैसे निकाला जाता है; और हम
इसे प्राकृतिक संख्याओं का अपना समुच्चय मानेंगे। अधिक सटीक रूप में:
\begin{defn}\ollabel{defnomega}
$I$ को अनंतता के अभिगृहीत द्वारा प्रदत्त कोई समुच्चय मानें। $s$ को फलन
$s(x) = x \cup \{x\}$ मानें। $\omega = \closureofunder{s}{\emptyset}$
रखें। हम $\omega$ के सदस्यों को \emph{प्राकृतिक संख्याएँ} कहते हैं और कहते
हैं कि $n$, $\emptyset$ पर $s$ के $n$ बार अनुप्रयोग का परिणाम है।
\end{defn}

अब आप पीछे देखकर जाँच सकते हैं कि कुछ अनुच्छेद पहले ``$n$'' अंकित समुच्चय
को संख्या $n$ \emph{माना} जाएगा।

इस नियतन के महत्त्व पर हम \olref[sth][z][nat]{sec} में चर्चा करेंगे। अभी
यह हमें एक सहज परिणाम सिद्ध करने देता है:

\begin{prop}\ollabel{naturalnumbersarentinfinite}
कोई भी प्राकृतिक संख्या डेडेकिंड-अनंत नहीं है।
\end{prop}

\begin{proof}
प्रमाण आगमन द्वारा है, अर्थात
\olref[sfr][infinite][induction]{thm:dedinfiniteinduction}। स्पष्टतः
$0 = \emptyset$ डेडेकिंड-अनंत नहीं है। आगमन चरण के लिए हम प्रतिधनात्मक
स्थापित करेंगे: यदि (असंगत रूप से) $s(n)$ डेडेकिंड-अनंत है, तो $n$
डेडेकिंड-अनंत है।

अतः मान लें कि $s(n)$ डेडेकिंड-अनंत है; अर्थात कोई !!{injection} $f$ ऐसा
है कि $\ran{f}\subsetneq \dom{f} = s(n) = n \cup \{n\}$। दो मामलों पर
विचार करना है।

\emph{मामला 1: $n \notin \ran{f}$।} अतः $\ran{f} \subseteq n$ और
$f(n) \in n$। $g = \funrestrictionto{f}{n}$ रखें; अब
$\ran{g} = \ran{f} \setminus \{f(n)\} \subsetneq n = \dom{g}$। इसलिए
$n$ डेडेकिंड-अनंत है।

\emph{मामला 2: $n \in \ran{f}$।}
$m \in \dom{f} \setminus \ran{f}$ को नियत करें और प्रांत
$s(n) = n \cup \{n\}$ वाला फलन $h$ इस प्रकार परिभाषित करें:
\[
h(x) = 
	\begin{cases}
		f(x) & \text{यदि }f(x) \neq n\\
		m & \text{यदि }f(x)=n
	\end{cases}
\]
अतः $h$ और $f$ हर जगह समान हैं, सिवाय इसके कि
$h(f^{-1}(n)) = m \neq n = f(f^{-1}(n))$। चूँकि $f$ कोई !!a{injection}
है, $n \notin \ran{h}$; और $\ran{h} \subsetneq \dom{h} = s(n)$। अब
मामला 1 के तर्क से $n$ डेडेकिंड-अनंत है।
\end{proof}

फिर भी यह प्रश्न बचता है कि अनंतता के अभिगृहीत का \emph{औचित्य} कैसे दिया
जाए। संक्षिप्त उत्तर है कि अपनी अब तक की कथा में एक और सिद्धांत जोड़ना
होगा। वह सिद्धांत यह है:
\begin{enumerate}
	\item[] \stagesinf. कोई अनंत चरण है। अर्थात ऐसा चरण है जो (a) पहला चरण
	नहीं है, (b) जिसके पहले कुछ चरण हैं, किंतु (c) जिसका कोई तात्कालिक
	पूर्ववर्ती नहीं है।
\end{enumerate}
अनंतता का अभिगृहीत इस सिद्धांत से सीधे निष्पन्न होता है। हम जानते हैं कि
प्राकृतिक संख्या $n$, चरण $n$ पर बनती है। अतः समुच्चय $\omega$ पहले अनंत
चरण पर बनता है। और $\omega$ स्वयं अनंतता के अभिगृहीत का साक्षी है।

किंतु इससे हम केवल इस प्रश्न पर लौट आते हैं कि \stagesinf{} का औचित्य कैसे
दिया जाए। \stagessucc{} की तरह यह संचयी-पुनरावर्ती पदानुक्रम की हमारी कथा
का स्पष्ट भाग नहीं था। इससे भी अधिक: पुनरावर्ती पदानुक्रम के मूल विचार में,
जहाँ समुच्चय चरण-दर-चरण बनते हैं, कुछ भी हमें यह मानने के लिए बाध्य नहीं
करता कि प्रक्रिया में कोई \emph{अनंत} चरण है। यह मानना पूरी तरह संगत लगता
है कि चरण प्राकृतिक संख्याओं की तरह क्रमित हैं।

पर इससे एक स्पष्ट समस्या उत्पन्न होती है।
\olref[sfr][infinite][dedekindsproof]{sec} में हमने डेडेकिंड के इस ``प्रमाण''
पर विचार किया था कि (विचारों का) कोई डेडेकिंड-अनंत समुच्चय है। संभव है वह
आपको बहुत संतोषजनक न लगा हो। किंतु यदि समुच्चय की पुनरावर्ती अवधारणा (अथवा
``विचार के नियम'') \stagesinf{} को हम पर ``अनिवार्य'' नहीं करती, तो किसी
डेडेकिंड-अनंत समुच्चय के अस्तित्व के दावे का अंतर्निहित औचित्य अब भी हमारे
पास नहीं है।

निश्चय ही यहाँ और बहुत कुछ कहा जा सकता है। किंतु आशा है कि अब आप इस पर
विचार आरंभ कर सकते हैं कि किसी अभिगृहीत (अथवा सिद्धांत) का औचित्य देने के
लिए \emph{क्या चाहिए}। आगे हम \stagesinf{} को बस मान लेंगे।

\end{document}
